\clearpage
\section*{KIRISH}
\addcontentsline{toc}{section}{KIRISH}

Bu bo’limdagi matn shablonning to’ldirgichidir. Mavzuning amaliy ahamiyati,
o’rganiladigan hodisa va tanlangan modelning qo’llanish chegaralarini shu yerda
ikki yoki uch xatboshi bilan bayon qiling. Haqiqiy manba tanlangach, uning
tegishli betiga tayangan fikrga havola bering~\cite{manba:nazariya}.

Nazariy tushunchalarni amaliy masalalar bilan bog’lang. Hisoblashlarda
qo’llanadigan farazlar, kattaliklarning birliklari va natijani tekshirish usuli
izchil tushuntiriladi. Amaliy to’plam ma’lumotlarini to’ldirgach,
unga murojaat qiling~\cite{manba:toplam}.

\textbf{Mustaqil ishning maqsadi} --- [mavzu] bo’yicha nazariy bilimlarni
tizimlashtirish va ularni masalalar yechishga qo’llash.

\textbf{Mustaqil ishning vazifalari:}
\begin{enumerate}
  \item mavzuga oid asosiy tushunchalar va farazlarni tushuntirish;
  \item zarur formulalarni keltirib chiqarish va birliklarini tekshirish;
  \item hisoblash natijalarini grafik va jadvallar yordamida tahlil qilish;
  \item o’quv qo’llanmadagi [N] mavzu bo’yicha mustaqil ishlash uchun
    berilgan 15 ta masalani yechish.
\end{enumerate}

Shablonda tuzilmani ko’rsatish uchun faqat ikkita sodda masala berilgan.
Ularni tanlangan mavzudagi masalalar bilan almashtiring va ro’yxatni to’ldiring.
